\section*{Chapter \ref{ch:m2:government-bonds} --- Government Bonds}

\begin{solution}{exo:m2:government-bonds:1}
Forty-one days of 184 have elapsed: $A = 2.125 \times 41/184 = 0.473505$ per
100.
\end{solution}

\begin{solution}{exo:m2:government-bonds:2}
$99 + 16.5/32 = 99.515625$. $101.296875 = 101 + 9.5/32$: 101-09+.
\end{solution}

\begin{solution}{exo:m2:government-bonds:3}
At 4\%: 100 (\cref{prop:m2:government-bonds:priceyield}). At 5\%: with $v =
1/1.025$, $2\sum_{k=1}^{6} v^k + 100 v^6 = 97.2459$.
\end{solution}

\begin{solution}{exo:m2:government-bonds:4}
Flows of 2, 2, 2 and 102 at 1, 2, 3, 4 half-years, $v = 1/1.02$: \omterm{def:m2:government-bonds:duration}{Macaulay
duration} 1.9419 years, modified $1.9419/1.02 = 1.9039$, \omterm{def:m2:government-bonds:duration}{DV01} $100 \times
1.9039 \times 10^{-4} = 0.019039$ per 100, USD~190.39 per million.
\end{solution}

\begin{solution}{exo:m2:government-bonds:5}
Duration alone: $+100.870911 \times 7.974966 \times 0.01 = +8.0444$. With
\omterm{def:m2:government-bonds:duration}{convexity}: $+8.0444 + \frac12 \times 100.870911 \times 75.824 \times 0.0001 =
+8.4268$. Full reprice at 3.20\%: $+8.4403$. For the fall, as for the rise, the
second-order estimate is within 0.014 of the truth; the first-order one misses
by 0.40.
\end{solution}

\begin{solution}{exo:m2:government-bonds:6}
Actual/actual: 44 days of a 181-day period, $2.5 \times 44/181 = 0.607735$.
30/360: from 15 February to 31 March counts $30 + 16 = 46$ days, $5 \times 46/360
= 0.638889$. The difference, 0.031 per 100 (USD~312 per million), matters to
anyone who holds a bond of one convention hedged with, or financed against,
an instrument of the other: every bond's terms name its day count, and a
system that assumes one for all will misprice the others.
\end{solution}

\begin{solution}{exo:m2:government-bonds:7}
Ten-year discount factor 0.657651, zero-coupon rate 4.2350\% against a \omterm{def:m2:government-bonds:zero}{par
yield} of 4.21\%. The \omterm{def:m2:government-bonds:coupon}{par} curve rises from two to ten years, so the earlier
\omterm{def:m2:government-bonds:coupon}{coupons} of a ten-year \omterm{def:m2:government-bonds:coupon}{par} bond are discounted at lower zero rates than its
final payment; its yield is a weighted average of those rates, below the
ten-year zero rate.
\end{solution}

\begin{solution}{exo:m2:government-bonds:8}
Only a zero-coupon bond has a duration equal to its life. A 4.5\% thirty-year
bond at \omterm{def:m2:government-bonds:coupon}{par} has a \omterm{def:m2:government-bonds:duration}{Macaulay duration} of 16.74 years and a \omterm{def:m2:government-bonds:duration}{modified duration} of
16.37; a one-basis-point rise costs $16.37 \times 0.0001 = 0.164\%$ of its price,
USD~1\,637 per million, about half of the claim.
\end{solution}

\begin{solution}{pb:m2:government-bonds:1}
\textbf{1.} The maturity, 30 September, is a month-end, so the \omterm{def:m2:government-bonds:coupon}{coupon} dates
are 31 March and 30 September: previous 31 March 2026, next 30 September 2026,
five days after settlement.
\textbf{2.} $A = 2 \times 178/183 = 1.945355$; \omterm{def:m2:government-bonds:accrued}{clean price} 100.224893, quoted
100-07 to the nearest 64th.
\textbf{3.} Dirty 102.170249: USD~102\,170\,249.
\textbf{4.} Five-year USD~451.43 per million; ten-year USD~804.44 per million.
\textbf{5.} 23.18 and 75.82.
\textbf{6.} $100 \times 451.43/804.44 = \text{USD}~56.12$ million of the ten-year
(ratio 0.5612).
\textbf{7.} The short leg's dirty value is $56\,117\,663 \times 1.00870911 =
\text{USD}~56\,606\,398$; net $102\,170\,249 - 56\,606\,398 = \text{USD}~45\,563\,851$
long, to be financed (\cref{ch:m2:repo-and-specials}).
\textbf{8.} $-\text{USD}~9.61$: the position has no \omterm{def:m2:government-bonds:duration}{DV01}, only a trace of
\omterm{def:m2:government-bonds:duration}{convexity}.
\textbf{9.} $+25$: $-\text{USD}~5\,934$. $-25$: $-\text{USD}~6\,092$.
\textbf{10.} The position is short \omterm{def:m2:government-bonds:duration}{convexity}: per unit of \omterm{def:m2:government-bonds:duration}{DV01}, the ten-year
has more \omterm{def:m2:government-bonds:duration}{convexity} ($75.8/7.97$ against $23.2/4.42$), and the desk is short
it. Whatever the direction, the short leg's price moves by more than its
duration predicts in the unfavourable way.
\textbf{11.} $-\text{USD}~450\,252$.
\textbf{12.} $+\text{USD}~449\,295$.
\textbf{13.} A rise of the ten-year yield relative to the five-year: the
position is a 5s10s \emph{steepener}.
\textbf{14.} Its \omterm{def:m2:government-bonds:duration}{DV01} is zero by construction, and so is its duration (\omterm{def:m2:government-bonds:duration}{DV01}
divided by value): the 45.6 million of net value carries no first-order rate
risk.
\textbf{15.} No. Time shortens both notes and the five-year's \omterm{def:m2:government-bonds:coupon}{coupon} of 30
September leaves its \omterm{def:m2:government-bonds:accrued}{dirty price}; after 28 days the long leg's \omterm{def:m2:government-bonds:duration}{DV01} is
USD~44\,519 and the short leg's USD~44\,863: net short USD~344 per basis point.
The hedge must be rebalanced.
\textbf{16.} A basis point moves a ten-year note almost twice as much as a
five-year: equal face or equal value would leave a large directional
position. \omterm{def:m2:government-bonds:duration}{DV01} matches the price change per unit of yield change.
\textbf{17.} Curve risk (the two yields move differently), \omterm{def:m2:government-bonds:duration}{convexity}, and the
drift of the \omterm{def:m2:government-bonds:duration}{DV01s}; also the financing of each leg in repo (\cref{ch:m2:repo-and-specials}).
\textbf{18.} Quotes and trade prices are in 32nds by convention; risk,
P\&L and hedge ratios are continuous quantities that must not be rounded to
a 64th before they are summed.
\textbf{19.} Named result: \emph{the DV01-neutral hedge} sells USD~56.12 million
of the ten-year against USD~100 million of the five-year (ratio 0.5612); a
25-basis-point parallel move either way costs about USD~6\,000
($-5\,934$ up, $-6\,092$ down), the price of being short \omterm{def:m2:government-bonds:duration}{convexity}, while a
10-basis-point steepening earns USD~449\,295.
\textbf{20.} On the shape of the curve, not its level.
\end{solution}

\begin{solution}{iq:m2:government-bonds:1}
The \omterm{def:m2:government-bonds:accrued}{dirty price} is what the buyer pays: the \omterm{def:m2:government-bonds:accrued}{clean price} plus the interest
accrued since the last \omterm{def:m2:government-bonds:coupon}{coupon}, which belongs to the seller. The \omterm{def:m2:government-bonds:accrued}{clean price}
removes the sawtooth of accrual and \omterm{def:m2:government-bonds:coupon}{coupon} payment, so its changes reflect
changes in yield only; that is why screens quote it and why yields are quoted
from it.
\iqlookfor{\omterm{def:m2:government-bonds:accrued}{accrued interest} as the difference, and why a quoted price must
not jump on \omterm{def:m2:government-bonds:coupon}{coupon} dates.}
\end{solution}

\begin{solution}{iq:m2:government-bonds:2}
The price is a sum of discount factors $(1 + y/f)^{-t}$, each convex in $y$.
\omterm{def:m2:government-bonds:duration}{Convexity} means the bond gains more for a fall in yield than it loses for an
equal rise, which is good for a holder, other things equal; markets price it,
so a position long \omterm{def:m2:government-bonds:duration}{convexity} usually gives up yield (carry) for it.
\iqlookfor{the sum of convex terms, and the carry-for-convexity trade-off.}
\end{solution}

\begin{solution}{iq:m2:government-bonds:3}
The zero: its duration is ten years, the \omterm{def:m2:government-bonds:coupon}{coupon} bond's is shorter (about
eight years at a 5\% yield), since earlier \omterm{def:m2:government-bonds:coupon}{coupons} pull the average time in.
Per 100 of face the comparison of \omterm{def:m2:government-bonds:duration}{DV01} needs the prices too: the zero trades
far below \omterm{def:m2:government-bonds:coupon}{par} (about 61 at 5\%), so its \omterm{def:m2:government-bonds:duration}{DV01} per 100 face ($61 \times 9.75 \times
10^{-4} \approx 0.059$) is smaller than the \omterm{def:m2:government-bonds:coupon}{coupon} bond's ($100 \times 7.8 \times
10^{-4} \approx 0.078$). Duration is per unit of value, \omterm{def:m2:government-bonds:duration}{DV01} per unit of face.
\iqlookfor{separating duration (relative) from \omterm{def:m2:government-bonds:duration}{DV01} (absolute).}
\end{solution}

\begin{solution}{iq:m2:government-bonds:4}
Sell the ten-year in face $N_{10} = N_5 \times \mathrm{DV01}_5/\mathrm{DV01}_{10}$,
roughly 0.55 to 0.6 of the five-year face at current yields. Remaining:
curve risk (you are now a steepener), \omterm{def:m2:government-bonds:duration}{convexity} (short), \omterm{def:m2:government-bonds:duration}{DV01} drift over time
and with yields, repo financing and specialness of each leg, and bid--ask on
rebalancing.
\iqlookfor{\omterm{def:m2:government-bonds:duration}{DV01} ratio, and naming the curve position that the hedge
creates.}
\end{solution}

\begin{solution}{iq:m2:government-bonds:5}
Month-end maturities (the end-of-month rule, including February in leap
years); settlement on a \omterm{def:m2:government-bonds:coupon}{coupon} date (zero accrued); settlement in an
ex-dividend period (negative accrued, for gilts); each market's day count
(actual/actual for Treasuries and gilts, 30/360 where a bond's terms say
so);
annual against semiannual \omterm{def:m2:government-bonds:coupon}{coupons}; short and long first \omterm{def:m2:government-bonds:coupon}{coupons}; the
regulation's and the exchange's published worked examples reproduced to the
printed digit; round trips between price and yield.
\iqlookfor{published worked examples as oracles, and the calendar edge
cases.}
\end{solution}

\begin{solution}{iq:m2:government-bonds:6}
A perpetuity paying $c$ a year prices at $P = c/y$. Then $dP/dy = -c/y^2$ and
the \omterm{def:m2:government-bonds:duration}{modified duration} is $-P'/P = 1/y = 25$ years at 4\%. Its \omterm{def:m2:government-bonds:duration}{Macaulay duration}
is $(1 + y)/y = 26$ years: the present-value weighted time of the payments,
$\sum_k k v^k / \sum_k v^k = 1/(1 - v) = (1+y)/y$.
\iqlookfor{the closed form and the difference between the two durations.}
\end{solution}
